% !TEX root = ../ICRA19Hybrid.tex

\section{RELATED WORK}
\label{sec:related_work}

\subsection{Hybrid Force-Velocity Control} % (fold)
\label{sub:hybrid_force_velocity_control}
The idea of using hybrid force-velocity control for manipulation under constraints can date back to 1980s.
Mason \cite{mason1981compliance} introduced a framework for identifying force and velocity controlled directions in a task frame given a task description. Raibert and Craig \cite{raibert1981hybrid} completed the framework and demonstrated a working system.
Yoshikawa \cite{yoshikawa1987dynamic} investigated hybrid force-velocity control in joint space under Cartesian space constraints, and proposed to use gradient of the constraints to find the normal of the constraint surface in the robot joint space.
There are also works on modeling the whole constrained robot system using Lagrange dynamics, such as analyzing the system stability under hybrid force-velocity control \cite{mcclamroch1988feedback}, or performing Cartesian space tracking for both positions and forces \cite{mills1989force}. Most of these works modeled only the robot and a rigid environment without any un-actuated degree-of-freedoms in the system. As an exception, Uchiyama and Dauchez performed hybrid force-velocity control for a particular example: two manipulators contacting one object \cite{uchiyama1988symmetric}.

There are lots of works on how to implement hybrid force-velocity controls on manipulators. For example, stiffness control can be used for this purpose. Velocity control is essentially a high stiffness control; force control can be implemented by low stiffness control with force offset.
Salisbury \cite{salisbury1980active} described how to perform stiffness control on arbitrary Cartesian axes with a torque-controlled robot.
Raibert and Craig \cite{raibert1981hybrid} divided Cartesian space into force/velocity controlled parts, then controlled them with separated controllers.
The impedance control \cite{hogan1985impedance}  and operational space control \cite{khatib1987unified} theory provided detailed analysis for regulating the force related behaviors of the end-effector for torque-controlled robots. Maples and Becker described how to use a robot with position controlled inner loop and a wrist-mounted force-torque sensor to do stiffness control on Cartesian axes \cite{maples1986experiments}. Lopes and Almeida enhanced the impedance control performance of industrial manipulators by mounting a high frequency 6DOF wrist \cite{lopes2008force}. Whitney \cite{whitney1987historical} and De Schutter \cite{de1998force} provided overviews and comparisons for a variety of force control methods.

\subsection{Motion Planning through Contacts} % (fold)
\label{sub:motion_planning_through_contacts}
Recently, a lot of works tried to solve manipulation under constraints without explicitly using force control. For holonomic constraints, De Schutter \etal proposed a constraint-based motion planning and state estimation framework \cite{de2007constraint}. Berenson \etal did motion planning on the reduced manifold of the constrained state space \cite{berenson2009manipulation}. For non-holonomic constraints, the most popular example is pushing \cite{lynch-mason-pushing,zhou2016convex,dogar2011framework}. Chavan-Dafle \etal performed in-hand manipulation by pushing the object against external contacts \cite{prehensile}. In these works, the robots interacted with the objects in a way that force control was not necessary.

